\documentclass[11pt]{amsart}
\usepackage{amsmath,amssymb,amsthm,enumitem}
\newtheorem{theorem}{Theorem}
\newtheorem{conjecture}[theorem]{Conjecture}
\newtheorem{proposition}[theorem]{Proposition}
\newtheorem{remark}[theorem]{Remark}
\title{R=PT Convergence: Shapovalov Eigenvalue Asymptotics\\
on the Virasoro Principal Series}
\begin{document}
\maketitle

\section{The problem}

The bar-cobar spectral $R$-matrix of the Virasoro algebra,
restricted to the principal series $\alpha = Q/2 + iP$,
is conjectured to equal the Ponsot--Teschner fusion kernel.
Eberhardt's uniqueness theorem (arXiv:2309.11540) reduces this
to a single analytic hypothesis: the bar-cobar $R$-matrix is
meromorphic in the external momenta~$\alpha_i$.

\begin{proposition}[Proved: level-by-level rationality]
At each fixed pair of Verma module levels $(N_1, N_2)$, the
matrix elements of $R^{\mathrm{bar},(N_1,N_2)}_{\alpha_1,\alpha_2}(z)$
are rational functions of $h_i = \alpha_i(Q - \alpha_i)$ with
poles exactly on the null-vector locus $h_i = h_{r,s}$.
\end{proposition}

The gap: does the infinite sum $\sum_N R^{(N)}$ converge to a
meromorphic function?  This reduces to bounds on the smallest
eigenvalue $\lambda_{\min}(G_N)$ of the Virasoro Shapovalov
(Gram) matrix at level~$N$.

\section{Numerical evidence}

At $c = 25$, $h = 2$ ($b = 1$, $Q = 2$, $P = 1$), exact computation
of the Virasoro Gram matrices gives:
\[
\begin{array}{c|cccccccc}
N & 1 & 2 & 3 & 4 & 5 & 6 & 7 & 8 \\
\hline
\lambda_{\min}(G_N) & 4 & 14.8 & 49.2 & 118.5 & 245.4 & 408.9 &
668.8 & 975.7 \\
p(N) & 1 & 2 & 3 & 5 & 7 & 11 & 15 & 22 \\
\lambda_{\min}/p(N) & 4.0 & 7.4 & 16.4 & 23.7 & 35.1 & 37.2 &
44.6 & 44.4
\end{array}
\]

$\lambda_{\min}$ grows monotonically and faster than $p(N)$.
The ratio $\lambda_{\min}/p(N)$ stabilises around~$44$.
This strongly suggests absolute convergence of the level sum on
the principal series.

\section{What rigorous proof requires}

\begin{enumerate}
\item A lower bound $\lambda_{\min}(G_N) \ge f(N)$ for some
  explicitly growing function~$f$, valid on the principal series
  at irrational~$b^2$.

\item The Kac determinant $\det G_N = C_N \prod_{rs \le N}
  (h - h_{r,s})^{p(N-rs)}$ gives the product of ALL eigenvalues,
  not the smallest.  The gap between $\det G_N$ (super-exponential
  growth) and $\lambda_{\min}$ (apparently polynomial growth)
  is the condition number $\kappa(G_N) = \lambda_{\max}/\lambda_{\min}$.

\item At $c = 25$, $h = 2$: $\lambda_{\max}$ grows like
  $\sim 10^{2.4 N}$ while $\lambda_{\min} \sim N^2$.  The condition
  number grows exponentially, but $\lambda_{\min}$ itself grows,
  which is what matters for convergence.

\item A uniform bound across compact subsets of the principal
  series would give meromorphicity via Weierstrass.
\end{enumerate}

\section{Attack paths}

\begin{enumerate}
\item \textbf{Explicit Shapovalov diagonalisation.}
  At each level~$N$, the Gram matrix $G_N$ is a $p(N) \times p(N)$
  matrix with entries polynomial in $h$ and~$c$.  The eigenvalues
  are algebraic functions of $h, c$.  On the principal series,
  $h = Q^2/4 + P^2$ is real and $\ge Q^2/4$.  All eigenvalues are
  positive (unitarity).  The smallest eigenvalue is bounded below
  by the minimum of a positive algebraic function on a compact set.

\item \textbf{Recursive structure.}
  By Cauchy interlacing, $\lambda_{\min}(G_{N+1}) \le
  \lambda_{\min}(G_N)$ does NOT hold (the matrices are not nested
  in the standard principal-submatrix sense).  Instead, $G_{N+1}$
  has $G_N$ as a block but with coupling terms from $L_{-1}$
  acting between levels.  The coupling terms are bounded by
  $\|L_{-n}\|_{\mathrm{op}} \sim n$ on level-$N$ states.

\item \textbf{Roussillon's series representation.}
  Roussillon (arXiv:2405.09325) gives a series representation of
  the fusion kernel at all irrational central charges.  If the
  bar-cobar $R$-matrix matches this series term by term (which
  follows from level-by-level rationality + matching at degenerate
  points), meromorphicity follows from Roussillon's convergence
  analysis.
\end{enumerate}

\end{document}
